\section*{Capítulo \ref{ch:b2:limb-organogenesis} --- Organogênese: a formação do membro dos tetrápodes}

\begin{solution}{exo:b2:limb-organogenesis:1}
Proximodistal: a \omterm{def:b2:limb-organogenesis:bud}{crista ectodérmica apical}, ao longo da ponta, FGFs.
Anteroposterior: a \omterm{def:b2:limb-organogenesis:bud}{zona de atividade polarizante}, na margem posterior,
\omterm{prop:b2:limb-organogenesis:zpa}{Sonic hedgehog}. Dorsoventral: o ectoderma dorsal, Wnt7a.
\end{solution}

\begin{solution}{exo:b2:limb-organogenesis:2}
Estilopódio: úmero, fêmur (Hoxa/d 9--10). Zeugopódio: rádio e ulna,
tíbia e fíbula (Hox 11). Autopódio: punho e mão, tornozelo e pé
(Hox 13).
\end{solution}

\begin{solution}{exo:b2:limb-organogenesis:3}
Campo do membro especificado (Tbx5); o broto cresce sob o FGF da
crista; a ZPA e o ectoderma dorsal o padronizam; as células que deixam
a \omterm{prop:b2:limb-organogenesis:aer}{zona de progresso} se condensam em raios; a cartilagem se diferencia;
as articulações se formam onde a cartilagem é suprimida; vasos
sanguíneos invadem e o osso substitui a cartilagem; as células
interdigitais morrem; os precursores musculares vindos dos somitos
migram para dentro, se fundem e se prendem por tendões.
\end{solution}

\begin{solution}{exo:b2:limb-organogenesis:4}
(a) Um coto com pouco ou nenhum úmero; (b) úmero, rádio e ulna, mas
nenhuma mão; (c) um membro completo --- o FGF substitui a crista.
\end{solution}

\begin{solution}{exo:b2:limb-organogenesis:5}
$x = 40\ln 2 = \qty{28}{\micro m}$, $40\ln 5 = \qty{64}{\micro m}$,
$40\ln 20 = \qty{120}{\micro m}$: territórios de 28, 36 e
\qty{56}{\micro m}.
\end{solution}

\begin{solution}{exo:b2:limb-organogenesis:6}
(a) 4-3-2-2-3-4, um conjunto especular completo. (b) Uma fonte mais
fraca dá um pico mais baixo, de modo que só o dedo de limiar baixo é
acrescentado: 2-2-3-4. (c) Uma exposição curta também só especifica o
dedo extra mais anterior: 2-2-3-4 --- as células integram a
concentração ao longo do tempo.
\end{solution}

\begin{solution}{exo:b2:limb-organogenesis:7}
$5000\times 2^{8} = 1.3\times 10^{6}$. Chegar a $4\times 10^{6}$ exige
$\log_2 800 = 9.6$ duplicações, de modo que o ciclo precisa durar em
média cerca de \qty{10}{h}, ou a contagem precisa incluir os
precursores musculares que migram dos somitos; a morte celular
agravaria a discrepância, em vez de reduzi-la.
\end{solution}

\begin{solution}{exo:b2:limb-organogenesis:8}
As células interdigitais da galinha recebem um sinal de BMP e morrem;
as células interdigitais do pato expressam um inibidor de BMP e
sobrevivem, de modo que a membrana persiste. O inibidor numa
microesfera salva a membrana da galinha: um pé palmado. A morte no
prazo previsto num enxerto significa que as células já tinham sido
instruídas antes da retirada --- a decisão é tomada cedo e executada
mais tarde, de forma autônoma.
\end{solution}

\begin{solution}{exo:b2:limb-organogenesis:9}
Sem Hoxa13 e Hoxd13: nenhum autopódio --- o membro anterior termina no
punho. Sem Hoxa11 e Hoxd11: um rádio e uma ulna muito encurtados, com a
mão presa quase diretamente ao úmero. Cada segmento precisa de seu
próprio par Hox; os domínios não são meros marcadores, mas instruções.
\end{solution}

\begin{solution}{exo:b2:limb-organogenesis:10}
O FGF da crista mantém o Shh ativo na ZPA; o Shh, por meio de um
revezamento no mesoderma, mantém o FGF ativo na crista. O sinal de
cada centro atesta, portanto, que o outro está funcionando, de modo que
um membro nunca cresce sem ser padronizado nem é padronizado sem
crescer, e uma perda transitória de um dos dois é reparada pelo outro.
Se o Shh não conseguisse manter o FGF, a crista regrediria em um dia e
o crescimento pararia no segmento que tivesse sido alcançado: um membro
truncado. A alça do parto é encerrada pelo evento que ela produz --- o
nascimento remove o estímulo; a alça do membro é encerrada pela
diferenciação --- o mesoderma para de fazer o revezamento quando o
membro está completo.
\end{solution}

\begin{solution}{exo:b2:limb-organogenesis:11}
$e^{4r} = 8$: $r = \ln 8/4 = \qty{0.52}{d^{-1}}$. Crescimento: uma
proliferação mais longa e mais rápida da cartilagem dos dedos (um
intensificador do membro que aumenta a expressão de um gene de
crescimento nos dedos). Morte: a sobrevivência do tecido interdigital
graças à expressão de um inibidor de BMP, que conserva a membrana.
\end{solution}

\begin{solution}{exo:b2:limb-organogenesis:12}
Os sinais são reutilizados porque um módulo receptor--transdução que
funciona é barato de reaproveitar: o que muda entre os tecidos não é o
sinal, mas o conjunto de genes que as células receptoras são
competentes para ativar, de modo que o mesmo perfil de concentração
significa neurônio motor num lugar e dedo 4 em outro. Uma \omterm{def:b2:genome-diversification:mutation}{mutação} num
gene desse tipo afeta, portanto, muitos órgãos ao mesmo tempo --- as
\omterm{def:b2:genome-diversification:mutation}{mutações} de Shh causam holoprosencefalia e defeitos dos membros juntos
---, a menos que ela esteja num intensificador específico de um tecido,
caso em que altera apenas um órgão, e é assim que ocorre a maior parte
das mudanças evolutivas nesses genes.
\end{solution}

\begin{solution}{pb:b2:limb-organogenesis:1}
\textbf{1.} $96/12 = 8$ duplicações: $2000\times 256 = 5.1\times 10^{5}$
células.
\textbf{2.} $1000/256 \approx 4$: o aumento das células e a matriz da
cartilagem.
\textbf{3.} Comprimento $\times 10$; um cilindro de raio constante
daria $\times 10$ em volume; o fator adicional de cem significa que o
raio também cresceu cerca de dez vezes --- o broto se alargou num remo à
medida que se alongava.
\textbf{4.} Aos 3 dias, $300/400 = \qty{75}{\%}$; aos 7 dias,
$300/4000 =
\qty{7.5}{\%}$.
\textbf{5.} Só a ponta está sob o sinal; à medida que a ponta avança,
as células deixadas para trás saem da zona na ordem em que foram
produzidas, de modo que as estruturas proximais são especificadas
primeiro e as distais por último.
\textbf{6.} O zeugopódio (rádio e ulna), especificado entre 3.5 e 4
dias.
\textbf{7.} Estilopódio até 3.5 dias, zeugopódio até 4.0 dias,
autopódio até 4.5 dias (as pontas dos dedos mais tarde).
\textbf{8.} Meio dia: um ciclo celular.
\textbf{9.} Cada ciclo passado na \omterm{prop:b2:limb-organogenesis:aer}{zona de progresso} faz avançar o
código Hox (9--10, depois 11, depois 13); as células que saem após um,
dois ou três ciclos levam o código do estilopódio, do zeugopódio ou do
autopódio.
\textbf{10.} Uma asa completa: o FGF é o que a crista fornecia.
\textbf{11.} Um segundo crescimento a partir da superfície dorsal:
estruturas distais duplicadas, um membro supranumerário.
\textbf{12.} O mesoderma para de retransmitir o Shh à crista, a
expressão de Shh termina, a crista regride e as células se diferenciam
em vez de se dividir: a alça que sustentava o crescimento se rompe.
\textbf{13.} $k = \ln 2/1740 = \qty{4.0e-4}{s^{-1}}$; $\lambda =
\sqrt{10^{-12}/4\times 10^{-4}} = \qty{50}{\micro m}$.
\textbf{14.} $x = 50\ln(1/0.6) = 26$, $50\ln 4 = 69$, $50\ln 12.5 =
\qty{126}{\micro m}$: territórios de 26, 43 e \qty{57}{\micro m}.
\textbf{15.} $600 - 126 = \qty{474}{\micro m}$.
\textbf{16.} Todas as fronteiras se deslocam para fora em
$50\ln 2 = \qty{35}{\micro
m}$: 61, 104, 161. As três larguras são 61, 43 e \qty{57}{\micro m}:
só o território do dedo 4 cresce, os demais apenas se deslocam. O padrão
avança para a frente em vez de ganhar um dedo, e a região anterior sem
dedos encolhe de 474 para \qty{439}{\micro m}.
\textbf{17.} Cada conjunto precisa de \qty{126}{\micro m}: pelo menos
\qty{252}{\micro m}; o broto de \qty{600}{\micro m} tem espaço, com
\qty{348}{\micro m} de tecido sem dedos no meio.
\textbf{18.} Com \qty{200}{\micro m}, os dois gradientes se sobrepõem e
se somam: as células do meio recebem o bastante para o dedo 3 e os
territórios do dedo 2 desaparecem --- 4-3-4 ou 4-3-3-4, menos dedos
que um conjunto especular completo.
\textbf{19.} Um quarto das células dá $c_0/4$: nenhuma célula atinge o
limiar do dedo 4 ou do 3 ($0.25 c_0$ na própria margem), e o limiar do
dedo 2 é ultrapassado até $50\ln(0.25/0.08) =
\qty{57}{\micro m}$: um único dedo 2 extra.
\textbf{20.} $3\times 300\times 300\times 200 = 5.4\times 10^{7}$
\unit{\micro m^3}: $5.4\times 10^{4}$ células; em \qty{86400}{s},
$0.6$ célula por segundo.
\textbf{21.} Microesfera de BMP na membrana do pato: a membrana morre e
os dedos se separam; inibidor na membrana da galinha: a membrana
sobrevive, um pé palmado.
\textbf{22.} $e^{0.5\times 4} = e^{2} = 7.4$. A segunda modificação é a
sobrevivência da membrana interdigital graças a um inibidor de BMP.
\textbf{23.} A fase tardia de expressão de Hox 13 que especifica um
autopódio, e a persistência da crista: a dobra do peixe forma raios de
nadadeira em vez disso. A nadadeira do \emph{Tiktaalik} contém ossos
semelhantes aos do punho --- uma nadadeira começando a ser um membro.
\textbf{24.} $5.4\times 10^{4}/1000 = 54$: $\log_2 54 = 5.8$
duplicações, cerca de \qty{69}{h}, três dias para produzir o que um dia
remove.
\textbf{25.} Oito duplicações; fronteiras dos dedos a 26, 69 e
\qty{126}{\micro m}; uma duplicação especular exige pelo menos
\qty{252}{\micro m}.
\end{solution}
